% main.tex
% this file contains the main typesetting information for the
% tetra text book.  it also contains front and back matter
% toc, etc. but pulls the chapters from individual files

\documentclass[14pt,twoside,openany]{memoir}
\usepackage{hyperref}
\usepackage{wallpaper}
\usepackage{tabularx}
\usepackage{textcomp}
\usepackage{graphicx}
\usepackage[labelfont=bf, font=bf, justification=raggedright,
    singlelinecheck=false, font=sf]{caption}
\usepackage{geometry}
\usepackage{color}
\usepackage{sectsty}
\usepackage{alltt}
\usepackage{float}
\usepackage{tcolorbox}

% define all the important meta information
\title{The Tetra Programming Language}
\author{Ian Finlayson, Ph.D. \\ Department of Computer Science \\
    The University of Mary Washington}
\date{\today}
\newcommand{\tpledition}{1\textsuperscript{st} Edition}

% fonts
\usepackage[T1]{fontenc}
\usepackage[bitstream-charter]{mathdesign}
\usepackage{FiraSans}
\usepackage{FiraMono}

% source code colors
\definecolor{keywordcolor}{RGB}{0, 64, 136}
\definecolor{typecolor}{RGB}{178, 140, 16}
\definecolor{builtincolor}{RGB}{115, 55, 140}
\definecolor{commentcolor}{RGB}{26, 93, 17}
\definecolor{valuecolor}{RGB}{170, 0, 0}
\definecolor{linenocolor}{RGB}{70, 70, 70}

% color of the background used for output boxes
\definecolor{outputbgcolor}{RGB}{245, 245, 245}
\definecolor{inputcolor}{RGB}{26, 93, 17}

% set my custom chapter style (based on "ger")
\chapterstyle{default}
\renewcommand*{\chapnumfont}{\normalfont\HUGE\bfseries\sffamily}
\renewcommand*{\chaptitlefont}{\normalfont\HUGE\bfseries\sffamily}
\renewcommand{\printchaptername}{\normalfont\HUGE\bfseries\sffamily Chapter} 
\renewcommand*{\chapterheadstart}{
\mbox{}\\\mbox{}\rule[0pt]{\textwidth}{1.4pt}\par}
\setlength{\midchapskip}{20pt}
\renewcommand*{\printchaptertitle}[1]{\chaptitlefont #1
\\
\vspace*{-8pt}
\\\mbox{}\rule[5pt]{\textwidth}{1.4pt}}

% set up the page geometry, normal 8.5 by 11 with margins
\geometry{
  body={6.5in, 8.5in},
  left=1.0in,
  top=1.25in
}

% an environment used for output boxes
\newenvironment{outputbox}
{\begin{tcolorbox}[colback=outputbgcolor, arc=0pt]
 \begin{alltt}}
{\end{alltt}
 \end{tcolorbox}}

% set the section headings in sans font
\setsecnumformat{\sffamily\csname the#1\endcsname\quad}
\setsecheadstyle{\Large\bfseries\sffamily}

% get rid of the garish link colors
\hypersetup {
    colorlinks = false,
    hidelinks = true
}

% set the memoir page style
\pagestyle{plain}

% title page stuff
\makeatletter
\def\maketitle{%
  \null
  \ThisCenterWallPaper{1.75}{skyway.jpg}
  \thispagestyle{empty}%
  \vfill
  {\color{white} \begin{center}\leavevmode
    \normalfont
      {\raggedright {\fontsize{44pt}{44pt}\selectfont \sffamily \textbf{\@title}\par}%
       \vskip 12pt
       {\Huge \sffamily \textbf{\tpledition}\par}}%
    \vskip 4in
      {\raggedleft \Huge \sffamily \textbf{\@author}\par}%
  \end{center}}%
  \vfill
  \null
  \cleardoublepage
  }
\makeatother

% begin the main document
\begin{document}

% don't put in silly blank pages
\let\cleardoublepage\clearpage

% plop in the title page, defined above
\maketitle

% do all the front matter junk
\frontmatter

% copyright etc.
{
    \setlength{\parindent}{0pt}
    \setlength{\parskip}{12pt}

    % put all the stuff at the bottom of the page
    \null
    \vfill

    \textcopyright \the\year\ by Ian Finlayson

    \href{https://creativecommons.org/licenses/by-sa/4.0/}{\includegraphics{cc}}

    This work is licensed under a
    \href{https://creativecommons.org/licenses/by-sa/4.0/}{Creative Commons
    Attribution-ShareAlike 4.0 International License}.  Electronic versions
    of this book may be found freely online at \url{http://tetra-lang.org/book}.

    Should you have any questions, suggestions or questions regarding this text,
    please feel free to email the author at \href{mailto:finlayson@umw.edu}{\texttt{finlayson@umw.edu}}

    The cover image is a photograph of the Sunshine Skyway Bridge which spans
    the Tampa Bay in Florida.  The image is attributed to Taildragon at \\
    \url{http://en.wikipedia.org}
}
\pagebreak

% make the table of contents
\tableofcontents

% and now put in the actual chapters
\mainmatter

\chapter{Computers \& Programs}


\section{What is a Computer?}






\section{What is a Program?}





\section{What is Tetra?}




\section{First Programs}

\def \codelabel {ch1hello}
The first program we will look at simply prints "Hello World!" to the screen.
The code for this program is given in Listing \ref{\codelabel{}}.
\input{chapter1/hello.ttr.tex}



\def \codelabel {ch1input}
The next program we will look at does input as well.
It can be seen in Listing \ref{\codelabel{}}.
\input{chapter1/input.ttr.tex}

\section{Summary}


\section{Exercises}


\chapter{Variables and Expressions}


\section{Using Variables}



\section{Data Types}


\section{Expressions}



\section{Operators}






\section{Summary}


\section{Exercises}


\chapter{Selection Statements}

Oftentimes, programs will need to perform different actions in different
scenarios.  For instance, a program might need to test input to see if
it is erroneous, or check which menu option a user selected.  Many
algorithms also rely on checking conditions to achieve their results.

If we continue the programs as recipes metaphor, we can consider recipes
that call for different baking temperatures depending on whether the pan
you are using is metal or glass, or even depending on whether you live in
a high altitude area.

In this chapter, we will learn how to write programs to test for conditions
and do different things based on whether they have occurred or not.

\section{If statements}

Tetra uses the keyword \texttt{if} to check if a condition has occurred.
For instance, if we want the user to enter a positive number, we can use
an \texttt{if} statement for this purpose:

\def \codelabel {ch3if1}
\def \codecaption{Using if to check input}
\input{chapter3/if1.ttr.tex}

The output of this program can be seen below:

\input{chapter3/if1.ttr.out}



\section{Conditions}


\section{Boolean Variables}


\section{Boolean Expressions}

\section{Else Statements}



\section{Elif Statements}


\section{Summary}


\section{Exercises}


\chapter{Looping}

\section{While Loops}


\section{For Loops}



\section{Using Ranges}


\section{Summary}


\section{Exercises}


\chapter{Parallelism}


\section{Parallel Statements}

\section{Parallel Pitfalls}

\section{Parallel For Loops}
\section{Summary}


\section{Exercises}


\chapter{Functions}

\section{Calling Functions}


\section{Writing New Functions}



\section{Passing Parameters}

\section{Summary}


\section{Exercises}


\chapter{List Processing}

\section{Using Lists}


\section{Looping on Lists}


\section{List Operations}

\section{Summary}


\section{Exercises}


\chapter{Synchronization}


\section{Lock Statements}


\section{Named Locks}


\section{Background Statements}


\section{Wait Statements}


\section{Summary}


\section{Exercises}


\chapter{More Data Types}

\section{Tuples}


\section{Dictionaries}


\section{Summary}


\section{Exercises}


\chapter{Object-Oriented Programming}

\section{Using Objects}


\section{Creating Classes}


\section{Initializing Objects}


\section{Using Methods}



\section{Summary}


\section{Exercises}








\end{document}

